% =============================================================
% BODY (EN)
% =============================================================

\section{Activity Solution}
\vspace{2mm}

The selected paper for this activity is \textit{Sentiment Analysis in Social Networks and Online Search Engines: A Financial Predictive Approach}. The paper analyzes how sentiment analysis, applied to social networks, online search engines, financial news, and other digital sources, can complement traditional financial data to improve the understanding and prediction of market behavior.

\vspace{2mm}

The article uses a systematic mapping of the literature as its research strategy to examine the current state of sentiment analysis techniques applied to the financial domain. Based on the studies analyzed, it is possible to examine the problem addressed, the solution approach considered, the main design decisions that can be identified, and the related quality attributes.

\vspace{2mm}

The following sections answer the four questions proposed in the activity based on the reading and analysis of the selected paper.


\subsection{What Was The Problem?}
\vspace{2mm}

The main problem identified in the paper is the difficulty of accurately predicting variations in financial markets. According to the article, these variations are not explained only by historical data or traditional economic indicators, but also by several external factors, such as global economic changes, news, investor behavior, and social perceptions expressed through digital platforms.

\vspace{2mm}

The paper explains that many financial operators consider historical performance useful for anticipating future returns. However, the work expands this perspective by arguing that user sentiment also plays an important role. In this sense, the problem is not only statistical or financial, but also informational: traditional models may be incomplete if they do not consider unstructured data from social networks, online searches, financial news, or forums.

\vspace{2mm}

The article also highlights that consumers and investors are increasingly difficult to predict. This creates the need to incorporate information sources capable of reflecting opinions, emotions, and trends in real time.

\vspace{2mm}

From a software design perspective, this problem can be interpreted as the need to build solutions capable of integrating heterogeneous information, processing textual data, and transforming it into useful information that can complement the models used for financial decision-making.


\subsection{What Solution Does The Paper Propose?}
\vspace{2mm}

The paper does not present an implemented software tool as a concrete solution. The \textbf{systematic mapping of the literature} is the research strategy used by the author to study the problem and analyze the current state of the field.

\vspace{2mm}

Based on the studies reviewed, the identified solution approach consists of \textbf{integrating sentiment analysis with financial predictive models}, complementing traditional financial data with unstructured information obtained from social networks, news, search engines, and other digital platforms.

\vspace{2mm}

Sentiment analysis makes it possible to automatically classify texts according to their polarity, for example, positive, negative, or neutral. This information can become an additional signal for interpreting market behavior and can be used together with financial data within predictive models.

\vspace{2mm}

The approach is based on the idea that social networks, web news, Google trends, and forum discussions can reflect collective perceptions about companies, assets, or economic contexts. Combining these sources with predictive models seeks to improve decision-making, anticipate market movements, and contribute to more timely risk management.

\vspace{2mm}

The paper also indicates that predictive performance can improve when multiple information sources are combined. Therefore, the approach does not depend on a single data source, but instead proposes complementing traditional financial indicators with textual information obtained from the digital environment.


\subsection{What Design Decisions Appear?}
\vspace{2mm}

The paper does not present the detailed design of a specific application or define a particular software architecture. However, based on the approaches analyzed, several relevant decisions can be identified from the perspective of designing a sentiment-analysis solution for financial prediction.

\vspace{2mm}

First, there is a decision to \textbf{integrate multiple information sources}. The approach considers traditional financial data together with information from social networks, web news, Google trends, and forum discussions. This decision seeks to enrich the available information and prevent predictions from depending exclusively on historical market behavior.

\vspace{2mm}

Second, the use of \textbf{natural language processing and sentiment analysis} can be identified. Information obtained from digital sources is mainly unstructured text and must therefore be processed to obtain interpretable information, such as positive, negative, or neutral sentiment polarity.

\vspace{2mm}

From a design perspective, a solution of this type would require mechanisms capable of obtaining, processing, classifying, and representing textual information before incorporating it into predictive models. The paper does not specify the concrete implementation of these mechanisms, so this aspect should be understood as a design need derived from the studied approach rather than as an architecture defined by the author.

\vspace{2mm}

Another relevant decision is the \textbf{integration of sentiment information with financial predictive models}. Sentiment analysis is not presented as an independent source for decision-making, but as complementary information that can be combined with financial indicators and data to improve the interpretation of market trends.

\vspace{2mm}

The approach also involves the use of \textbf{information close to real time}. Social networks, news, and other digital platforms make it possible to rapidly observe changes in opinions and emerging trends. A solution based on this approach must therefore consider the frequent incorporation of new information so that its results remain useful within a changing financial context.

\vspace{2mm}

The paper also discusses the possibility of developing \textbf{personalized and adaptive strategies}. Recommendations may be adjusted according to investor risk profiles, preferences, and behaviors, meaning that the design should consider different users and decision contexts.

\vspace{2mm}

Finally, the approach requires considering aspects such as \textbf{robustness against noise and language ambiguity}, \textbf{result transparency}, and \textbf{financial information protection}. The article acknowledges that predictive models are not infallible and must handle subjective information, volatile markets, and risks associated with the use of financial data.


\subsection{Which Quality Attributes Does It Try To Improve?}
\vspace{2mm}

The most evident quality attribute associated with the approach is \textbf{accuracy}. The paper analyzes how incorporating information from social networks, search engines, and financial news can complement traditional financial data and contribute to more accurate predictions of market trends.

\vspace{2mm}

Another relevant attribute is \textbf{timeliness}. Digital sources can provide information in real time, making it possible to detect changes in opinion, emerging trends, or public reactions before they are fully reflected in traditional financial indicators. This can support faster decisions by investors, analysts, and companies.

\vspace{2mm}

\textbf{Adaptability} is also important. The article explains that sentiment analysis can be used to personalize financial strategies according to investor risk profiles, preferences, and behaviors. Therefore, a solution based on this approach should be capable of adapting its recommendations to different users and decision contexts.

\vspace{2mm}

\textbf{Robustness} is another necessary attribute. Language used in social networks and news may be ambiguous, subjective, or contain irrelevant information, while financial markets are highly variable. A system of this type should be capable of distinguishing relevant information from noise and preventing unreliable data from producing weak conclusions.

\vspace{2mm}

\textbf{Security and privacy} are also relevant aspects. In its conclusions, the paper identifies the need to guarantee the privacy and security of financial data. This consideration is particularly important when systems work with information related to users, investors, or institutions.

\vspace{2mm}

Finally, the article highlights the importance of \textbf{transparency and understandability}. The complexity of predictive models must be balanced with the user's ability to understand the results obtained. Therefore, producing a prediction is not sufficient by itself: the results should also be presented clearly enough to support responsible decision-making.


\section{Conclusion}
\vspace{2mm}

In summary, the paper addresses the difficulty of anticipating financial market behavior in contexts characterized by multiple variables and high volatility. To study this problem, the author uses a systematic mapping of the literature and, based on the studies analyzed, identifies the integration of sentiment analysis with financial predictive models as a relevant solution approach.

\vspace{2mm}

The studied approach seeks to complement traditional financial data with unstructured information from social networks, search engines, news, and other digital platforms. In this way, sentiment expressed by users and investors can become an additional source of information for interpreting trends and supporting decision-making.

\vspace{2mm}

From a software design perspective, several decisions can be identified regarding the integration of multiple data sources, natural language processing, incorporation of sentiment information into predictive models, use of information close to real time, personalization, and the consideration of risks related to data quality and protection.

\vspace{2mm}

The main quality attributes associated with the approach are accuracy, timeliness, adaptability, robustness, security, privacy, transparency, and result understandability.

\vspace{2mm}

It is important to emphasize that the paper does not define a software architecture or the detailed design of a concrete solution. For this reason, the identified decisions correspond to design characteristics and needs that can be derived from the approaches studied, without attributing components, technologies, or patterns to the article that it does not explicitly specify.